\section{Selective Disclosure}

\subsection{Selective Disclosure in JWT-Based Verifiable Presentations}

Selective disclosure in JSON Web Token (JWT)-encoded verifiable credentials enables data minimisation while preserving authentication guarantees. The holder can reveal only a subset of the claims originally issued, while verifiers still validate authenticity and integrity. \\
This mechanism, currently specified in an IETF draft, relies on salted-hash–based commitments rather than zero-knowledge proofs.

\subsubsection*{JWT Structure and Analogy with Verifiable Credentials}
A JWT consists of three Base64URL-encoded components: header, payload, and signature.  
\begin{itemize}
    \item \textbf{Header}: token type and signature algorithm.
    \item \textbf{Payload}: registered, public, and private claims (e.g., \texttt{iss}, \texttt{sub}, \texttt{aud}, \texttt{iat}, \texttt{nbf}, \texttt{exp}, \texttt{jti}).
    \item \textbf{Signature}: integrity and authenticity protection using the issuer's private key.
\end{itemize}
This structure mirrors verifiable credentials: metadata (header), claims (payload), and issuer signature.

\subsubsection*{Selective Disclosure JWT (SD-JWT)}
Selective disclosure is achieved by replacing cleartext claims with salted digests:
\begin{itemize}
    \item The \textbf{issuer} places a hash of each selectively disclosable claim in the JWT payload.
    \item The \textbf{holder} later discloses only selected claims in plaintext (each disclosure includes: salt, claim name, claim value).
    \item The \textbf{verifier} recomputes the digest and checks it against the SD-JWT.
\end{itemize}
This approach supports minimal disclosure but is not zero-knowledge: plaintext claims must be transmitted at presentation time.

\subsubsection*{Disclosure Mechanism}
A disclosure has the form:
\[
    \text{Disclosure} = \{\text{salt}, \text{claim}, \text{value}\}
\]
and its digest is:
\[
    \text{digest} = \text{SHA-256}(\text{Base64}( \text{Disclosure}))
\]
The SD-JWT includes:
\begin{itemize}
    \item the issuer-signed JWT containing all salted digests;
    \item zero or more disclosures, appended in compact serialization using \texttt{\textasciitilde} separators;
    \item optional cleartext claims designated by the issuer.
\end{itemize}

\subsubsection*{Issuance and Presentation Flow}
\begin{enumerate}
    \item The issuer produces an SD-JWT containing salted digests and indicates which claims are disclosable.
    \item The holder receives the SD-JWT plus all disclosures.
    \item The holder selects which disclosures to reveal to each verifier.
    \item The verifier:
    \begin{itemize}
        \item verifies the issuer signature over the SD-JWT,
        \item recomputes digests for the disclosed claims,
        \item checks that each digest appears in the SD-JWT,
        \item validates the SD-JWT as a standard JWT.
    \end{itemize}
\end{enumerate}

\subsubsection*{Key Binding for Holder Consent}
To prevent replay or theft of SD-JWTs, the draft defines a \textbf{key-binding JWT}:
\begin{itemize}
    \item The issuer embeds the holder’s public key in the SD-JWT (claim \texttt{cnf}).
    \item The holder proves possession of the corresponding private key by signing a hash of the SD-JWT, along with a nonce and audience.
    \item The verifier checks that the signature corresponds to the public key contained in the SD-JWT.
\end{itemize}
This binds the SD-JWT to a specific holder and ensures freshness and audience restriction.


